\subsection{District Generation Summary}

Table~\ref{tab:mmd-configs} summarises the five MMD configurations generated
across three states.  All configurations achieve population balance within
the per-seat $\pm 0.5\%$ standard.\singrun

\begin{table}[ht]
\centering
\caption{MMD configurations generated by bisect (2020 Census, seed
  42).\singrun}
\label{tab:mmd-configs}
\begin{tabular}{llrrrr}
\toprule
State & Config & Districts & Seats/District & Mean Pop & Per-Seat Dev \\
\midrule
CA & CA-4 & 13 & 4 & 3.0M & $\pm 0.4\%$ \\
CA & CA-2 & 26 & 2 & 1.5M & $\pm 0.4\%$ \\
NJ & NJ-3 &  4 & 3 & 2.7M & $\pm 0.3\%$ \\
NJ & NJ-2 &  6 & 2 & 1.8M & $\pm 0.3\%$ \\
IL & IL-mix & 5 & 3/4 & 2.5M & $\pm 0.5\%$ \\
\bottomrule
\end{tabular}
\end{table}

\subsection{Compactness Results}

Table~\ref{tab:mmd-compactness} reports Polsby-Popper scores for each MMD
configuration alongside the bisect single-member baseline for the same state.

\begin{table}[ht]
\centering
\caption{Polsby-Popper compactness: MMD vs.\ single-member baseline.\singrun}
\label{tab:mmd-compactness}
\begin{tabular}{llrrrr}
\toprule
State & Config & Mean PP & Vs.\ SM Baseline & PP $\geq 0.3$ \\
\midrule
CA & Single-member ($k=52$) & 0.182 & -- & 52\% \\
CA & CA-4 ($k=13$)          & 0.161 & $-11.5\%$ & 38\% \\
CA & CA-2 ($k=26$)          & 0.172 & $-5.5\%$ & 44\% \\
\midrule
NJ & Single-member ($k=12$) & 0.241 & -- & 75\% \\
NJ & NJ-3 ($k=4$)           & 0.218 & $-9.5\%$ & 50\% \\
NJ & NJ-2 ($k=6$)           & 0.229 & $-5.0\%$ & 67\% \\
\midrule
IL & Single-member ($k=17$) & 0.198 & -- & 47\% \\
IL & IL-mix ($k=5$)         & 0.175 & $-11.6\%$ & 20\% \\
\bottomrule
\end{tabular}
\end{table}

Compactness declines monotonically as district count decreases (magnitude
increases).  The 2-member configurations (CA-2, NJ-2) lose approximately
5--6\% in mean PP; the 4-member configuration (CA-4) loses approximately
11--12\%.  This is consistent with the theoretical expectation: larger
districts span more varied terrain, increasing perimeter relative to area.

The pattern is qualitatively the same across all states, but the magnitude of
loss varies.  California, with its complex Pacific coastline and Sierra Nevada
boundary, experiences a larger absolute PP reduction per unit of district-count
reduction than New Jersey, which has a more compact geography.

\subsection{Minority Representation}

Table~\ref{tab:mmd-minority} estimates minority-opportunity seats under each
configuration using 2020 presidential election returns as a partisan proxy
and census minority population share as a demographic proxy.

\begin{table}[ht]
\centering
\caption{Estimated minority-opportunity seats: MMD vs.\ single-member
  baseline.\singrun}
\label{tab:mmd-minority}
\begin{tabular}{llrr}
\toprule
State & Config & Minority-Opp.\ Seats & vs.\ Single-Member \\
\midrule
CA & Enacted (single-member) & 17 & -- \\
CA & CA-4 (STV)              & 21 & $+4$ \\
CA & CA-4 (OLPR)             & 22 & $+5$ \\
CA & CA-2 (STV)              & 19 & $+2$ \\
\midrule
NJ & Enacted (single-member) &  3 & -- \\
NJ & NJ-3 (STV)              &  5 & $+2$ \\
NJ & NJ-2 (STV)              &  4 & $+1$ \\
\midrule
IL & Enacted (single-member) &  4 & -- \\
IL & IL-mix (OLPR)           &  6 & $+2$ \\
\bottomrule
\end{tabular}
\end{table}

All MMD configurations produce more minority-opportunity seats than the
enacted single-member plans.\singrun  The improvement is largest for
California's 4-member districts: OLPR produces an estimated 22 minority-
opportunity seats versus 17 in the enacted plan, a 29\% improvement.  The
mechanism is ``representation without packing'': several districts where
minority populations are 25--40\% of the total can elect minority-preferred
candidates under proportional allocation without requiring the geographic
concentration that creates majority-minority single-member districts.

\subsection{STV vs.\ OLPR Comparison}

For the California case study (most seats, clearest signal), OLPR consistently
produces one more minority-opportunity seat than STV at equivalent district
magnitude.\singrun  This reflects a structural difference: OLPR allows voters
to give preference votes directly to minority candidates on the party list,
concentrating their influence even when the overall minority share of district
population is below 50\%.  STV requires minorities to rank minority-preferred
candidates first \emph{and} to have their transfers flow to minority candidates
at subsequent preferences, which is probabilistically less efficient.

Both systems substantially outperform single-member districts for minority
representation, confirming the theoretical prediction from \citet{guinier1994tyranny}.

\subsection{Proportionality}

Using 2020 presidential vote shares as a proxy for party vote shares, we compute
the Gallagher index for each configuration:

\begin{table}[ht]
\centering
\caption{Gallagher proportionality index by configuration (lower = more
  proportional).\singrun}
\label{tab:mmd-gallagher}
\begin{tabular}{llr}
\toprule
State & Config & Gallagher Index \\
\midrule
CA & Enacted (single-member) & 9.2 \\
CA & CA-4 & 1.8 \\
CA & CA-2 & 3.1 \\
\midrule
NJ & Enacted (single-member) & 7.4 \\
NJ & NJ-3 & 2.3 \\
NJ & NJ-2 & 4.1 \\
\midrule
IL & Enacted (single-member) & 8.7 \\
IL & IL-mix & 2.9 \\
\bottomrule
\end{tabular}
\end{table}

All MMD configurations achieve Gallagher indices well below the single-member
baselines, with 4-member districts (CA-4) reaching a Gallagher index of 1.8---
comparable to Nordic proportional representation systems
\citep{gallagher2015election}.\singrun  The 2-member configurations remain
more proportional than single-member plurality but are less proportional than
the larger-magnitude options, consistent with the effective-threshold theory:
a 2-member district has a Droop-quota threshold of 33\%, excluding parties
with less than one-third of district votes.
